\documentclass[11pt,a4paper]{article}
\usepackage[utf8]{inputenc}
\usepackage[british]{babel}
\usepackage[a4paper,margin=2.3cm]{geometry}
\usepackage{booktabs,longtable,array}
\usepackage[most]{tcolorbox}
\usepackage[hidelinks]{hyperref}
\setlength{\parindent}{0pt}
\setlength{\parskip}{4pt}
\newtcolorbox{note}[1]{colback=black!4,colframe=black!45,boxrule=0.6pt,arc=2pt,
  fonttitle=\bfseries\small,title={#1},fontupper=\small}

\title{Women in \textit{The Ladies' Diary}, 1704--1840\\[2pt]
\large What Leybourn's four volumes and the later Diaries show}
\author{Working notes for the booklet series}
\date{October 2026, revised}

\begin{document}
\maketitle

\section*{The short version}

\textit{The Ladies' Diary} was founded in 1704 as an almanac for women. In the first
twenty years women set about one question in ten. After 1724 they almost vanish as
setters. Through the rest of the century a handful of women's names appear, never more
than a few in a decade. By the early nineteenth century there is one.

As solvers they last longer. A small group of women answered regularly in the 1710s.
Then there are very few until the 1790s, when Miss Maria Middleton and others reappear.
In the 1800s Miss Mary Weston of Gainsborough answers eight questions.

After 1816, where Leybourn stops, the Diary's own pages show three more women setting
questions, and a scattering of solvers: Sarah Ellis, who went on answering as Mrs Sarah
Clay after her marriage; Jane Thoburn of Blyth, aged twelve; and Miss L.~L., who by 1834
was among the Diary's most regular solvers. In the last issue, for 1840,
no woman appears in the mathematics.

One correction comes out of this work. Leybourn's index identifies ``Miss Ann
Nicholls'', the most prolific woman setter of the middle of the century, as a pen-name
of William Wales. Wales was the astronomer who later sailed with Cook. Book Three originally
treated her as a woman. It has been corrected.

\section*{Sources and method}

For 1704 to 1816 everything here comes from Thomas Leybourn, \textit{The Mathematical
Questions, Proposed in the Ladies' Diary, and their Original Answers} (4 vols, London,
1817). We worked from the Internet Archive's machine-read text of all four volumes, and
from a machine-read copy of the Diary for 1780. Leybourn's volumes stop in 1816. For 1816
to 1840 we worked from three bound volumes of the Diary itself: the issues for 1816--1825
and 1825--1834, digitised by Google Books, and for 1835--1840, the Harvard College Library
copy digitised by HathiTrust.\footnote{\url{https://hdl.handle.net/2027/hvd.hxvmnt}.}
There is no index for these years, so the names after 1816 come from reading the pages
themselves, including the ``also answered'' lists, and they are less complete than
Leybourn's for the solvers he printed. The main tool is Leybourn's \emph{Index of the
Names of the Persons who have Proposed and Answered the Questions}, at the end of
volume~IV. It gives, for every contributor, the questions he or she set and answered.
Where the index was hard to read, we checked names against the text. The section
``What we know of them'' goes beyond Leybourn, and its sources are given in footnotes.

Three cautions.
\begin{enumerate}
\item \textbf{Pen-names.} Many contributors used pseudonyms. A woman's pen-name, such as
Amanda, Rosamond or Silvia, does not prove a woman. ``Ann Nicholls'' shows that men did
use them. Leybourn names the person behind some pseudonyms, but not all.
\item \textbf{The index covers only the printed answers.} Leybourn's index lists the
solvers whose answers he printed or named. The Diary itself also printed long lists of
other people who ``also answered''. Those are not in the index. So the index undercounts
women solvers. The 1780 issue, for example, names Mrs Eleanor Suggett as a solver of
three questions, and she is not in the index.
\item \textbf{Only the mathematics.} Leybourn collected only the mathematical questions.
The Diary's enigmas, rebuses, queries and poems, where women were far more numerous,
are not in his volumes at all. Even in 1840, when no woman appears in the mathematics,
about a fifth of the contributors to the enigmas were women.\footnote{J.~Albree and S.~H.~Brown, ``A valuable monument of mathematical genius'', \textit{Historia Mathematica} 36 (2009), 10--47, p.~32.}
\end{enumerate}

\section*{How many questions women set}

Counting questions set under a woman's name or a woman's pen-name, prize questions
included, and leaving out ``Ann Nicholls'':

\begin{center}
\begin{tabular}{l l r r r}
\toprule
Volume & Years & Questions & By women & Rate\\
\midrule
I   & 1704--1748 & about 330 & 11 & about 1 in 30\\
\quad of which 1709--1724 & & about 110 & 11 & about 1 in 10\\
\quad of which 1725--1748 & & about 220 & 0 & none\\
II  & 1749--1776 & about 430 & 7  & about 1 in 60\\
III & 1777--1802 & about 415 & 5  & about 1 in 80\\
IV  & 1802--1816 & about 225 & 1  & about 1 in 220\\
\textit{The Diary} & 1817--1840 & about 375 & 4  & about 1 in 95\\
\bottomrule
\end{tabular}
\end{center}

Teri Perl made an independent count in 1979, of questions \emph{and} answers by
women, not only questions set. Albree and Brown summarise it: 42 contributions from
1710 to 1725, then a gap of 23 years, then 17 from 1748 to 1768, then only two in the
next 27 years.\footnote{J.~Albree and S.~H.~Brown, ``A valuable monument of
mathematical genius'', \textit{Historia Mathematica} 36 (2009), p.~17, from T.~Perl,
\textit{Historia Mathematica} 6 (1979).} The shape is the same as ours: a strong start,
a long silence, a short revival around 1750. The last row of our table
counts four questions by three women: Mary Knight and Susan Jackson in 1817, and Miss
L.~L. in 1834 and 1836. Albree and Brown also count three women setters between 1817 and
1836. Perl's figure for 1748--1768 probably
includes the ``Ann Nicholls'' questions, which we leave out.

\begin{note}{Corrections made to the booklets}
Book Three first said that ``at least seventeen'' of volume~II's questions carry a
woman's name, about one in twenty-four, ten of them by Ann Nicholls. It now says seven,
about one in sixty, and names Wales. Its three ``Nicholls'' problems carry both names.
The footnotes in Books Three and Four first gave ``about one in ten'' for volume~I as a
whole. They now give about one in thirty, with one in ten for 1709--1724.
\end{note}

\section*{Women who set questions}

\begin{longtable}{p{4.2cm} p{2.8cm} p{8.2cm}}
\toprule
Name as printed & Questions & Notes\\
\midrule
\endhead
\multicolumn{3}{l}{\textit{Volume I}}\\
Lydia Fisher & Q18 (c.\ 1711) & Book One, ``Lydia Fisher's numbers''.\\
Sarah Brown & Q29 (c.\ 1713) & Book One, ``The patient lover''.\\
Mrs Anna Wright & prize question 1713 & The star nearest the pole at the Creation (astronomy).
  She also answered Q23 and the prize question for 1711.\\
Mrs Barbara Sidway & Q36 (c.\ 1714) & The greatest cylinder in a cone: Book One, Purple.
  She also answered Q17, 20, 21, 22 and two prize questions.\\
Mrs Boydell & Q43 (c.\ 1715) & Book One, ``The old man's name''.\\
Xantippe, of Helmsley & Q55 (c.\ 1717) & Pen-name. Book One, ``The legacy''.\\
The Lady Mer.\ Heyshot & Q58 (c.\ 1717) & Book One, ``Two sons''.\\
Mrs Mary Nelson, of Newmarket & Q72 (c.\ 1717) & Book One, ``The prize money''. She also
  answered several questions and the prize questions for 1714 and 1715.\\
Adrastea & Q89 (c.\ 1721), prize questions 1722 and 1724 & Pen-name. Book One, ``The
  passengers' fares''. In one of her questions she drops stones into Eldon Hole, in Derbyshire, in 1718.\\
S.~T. & Q80 (c.\ 1720) & Gender not stated. Not counted.\\[4pt]
\multicolumn{3}{l}{\textit{Volume II}}\\
Amanda & Q296 (1748/9) & Pen-name. Book Three, ``Always one short''.\\
Miss Munlove (index: Manlove) & Q298 (1748/9) & Book Three, ``Placing the servants'' (and
  the grey box).\\
Rosamond & Q305 (1748/9) & Pen-name. Amicable numbers. Book Three, ``Seven more than
  itself''.\\
Miss Maria Atkinson & Q377 (1754), Q391 (1755) & Signed ``Miss Maria A---i---s---n'' for
  Q377. Q391 is her age in verse, which gives Book Three its epigraph.\\
Miss T.~S---e & Q434 (1758) & Book Three, ``Height, years and fortune''.\\
Miss S.~T. & Q463 (1760) & ``Miss'' in the text; Leybourn's index has only ``T.~S.''.
  Book Three, ``The drover's sheep''.\\[4pt]
\multicolumn{3}{l}{\textit{Volume III}}\\
Miss Nancy Mason, of Clapham, Yorkshire & Q984 (1795), Q1003 (1796), Q1018 (1797) & Her age (Book
  Four), the five sons' square field (Book Four), a silver vessel of wine.\\
Miss Sarah Cowen & Q1030 (1798), Q1059 (1800) & A field from a lost figure; tangent three
  times sine (Book Four).\\[4pt]
\multicolumn{3}{l}{\textit{Volume IV}}\\
Mrs Macdonald, of Spindlestone & Q1098 (1802) & The falling ball of cork. Leybourn's index
  has ``Macdonal, J.'' (\textit{Beyond A-level}).\\[4pt]
\multicolumn{3}{l}{\textit{The Diaries for 1816--1825 (not in Leybourn)}}\\
Miss Mary Knight, of Boston & Q1314 (1817) & Cone with slant 10; she answered it herself
  (\textit{After Leybourn}).\\
Miss Susan Jackson, of Mile End & Q1316 (1817) & Two balls from a tower
  (\textit{After Leybourn}).\\[4pt]
\multicolumn{3}{l}{\textit{The Diaries for 1826--1834 (not in Leybourn)}}\\
Miss L.~L. & Q1577 (1834) & The largest circle in a triangle cut from a semicircle. Her
  own answer, printed in 1835, is exactly right (\textit{The Gregory years}).\\[4pt]
\multicolumn{3}{l}{\textit{The Diaries for 1835--1840 (not in Leybourn)}}\\
Miss L.~L. & Q1608 (1836) & Also Q1577 above. Q1608 is the shortest route from London to
  Constantinople touching the equator (\textit{The last years}).\\
\bottomrule
\end{longtable}

\section*{Women who answered}

Up to 1816, taken from Leybourn's index; the Diary's own ``also answered by'' lists,
which are not indexed, would add more. After 1816, taken from the Diary's pages, lists
included.

\begin{longtable}{p{4.2cm} p{3.8cm} p{7.2cm}}
\toprule
Name as printed & Answers & Notes\\
\midrule
\endhead
\multicolumn{3}{l}{\textit{Volume I}}\\
A Lady & Q6 & ``Solution by a Lady'', the earliest in the volume.\\
Mrs Mary Wright & Q16, the prize question set in 1710 & Her answer printed in 1711, when she won the twelve Diaries. Book One, introduction.\\
Mrs Barbara Sidway & Q17, 20, 21, 22, two prize questions & Also a setter.\\
Mrs Anna Wright & Q23, prize 1711 & Also a setter.\\
Mrs Mary Nelson & several, prize questions 1714, 1715 & Also a setter.\\
Mrs Morgan & prize 1714 &\\
Silvia & Q41, 44, 45, 46 & Pen-name.\\
Hesychia & Q49, 67 & Pen-name.\\
Mrs Adway & Q56 &\\
Anna Philomathes & Q59, 62, 63, 64 & Pen-name.\\
Sylvia & Q93, 107 & Pen-name, perhaps the same person as Silvia.\\
Eliz.\ Dod & one question, about Q95 & The number is hard to read in the index.\\
Adrastea & prize 1722 & Also a setter.\\
A Lady, Grubean & Q170 &\\[4pt]
\multicolumn{3}{l}{\textit{Volume II}}\\
Miss E.~S. & Q391 & She answered Miss Atkinson's age question.\\
Amaryllis & Q462 & Pen-name. The question was ``Ann Nicholls''s, that is, William
  Wales's.\\[4pt]
\multicolumn{3}{l}{\textit{The Diary for 1780 (not in the index)}}\\
Mrs Eleanor Suggett & Q742, 743, 746 & In the ``also answered'' lists.\\
Mrs Jane Roberts & prize enigma & Won ten Diaries; set and answered rebuses.\\
Miss Diana Browne & prize enigma & Won eight Diaries.\\
Miss Nancy --- & rebuses & Answered in verse.\\[4pt]
\multicolumn{3}{l}{\textit{Volume III}}\\
Miss Betty Claxton, of Benwell & Q909 & First printed answer to the window and the wall
  (Book Four).\\
Miss Maria Middleton, of Eden, near Durham & Q1002, 1014, 1017, 1041 & Including the
  three heights of a triangle (Book Four).\\
Miss M.~Castieau, of Shropshire & Q1059 & Answered Miss Cowen's question.\\[4pt]
\multicolumn{3}{l}{\textit{Volume IV}}\\
Miss Mary Weston, of Gainsborough & Q1195, 1210, 1231, 1247, 1258, 1286, 1291, 1294 &
  The most persistent woman solver in the collection after 1720.\\
Miss Susan (Susannah) Jackson, of Mile End & Q1241, 1254, 1269 and one more &\\
Mrs Elizabeth Bryhon, of Helmsley & Q1157 &\\[4pt]
\multicolumn{3}{l}{\textit{The Diaries for 1816--1825 (not in Leybourn)}}\\
Miss Mary Knight & Q1299 (her solution printed, 1817), Q1314, and several in 1818 & \\
Miss Susan Jackson & several, 1817--1818 & Also a solver in Leybourn's vol.~IV.\\
Miss Mary R---e; Miss Mary West & 1818 lists & ``West'' may be Mary Weston.\\
Miss Clara Clay & Q1391 (1823) &\\
Miss Sarah Ellis, later Mrs Sarah Clay & Q1406 (1824); Q1423, Q1429 (1825) & Printed in
  1825 as ``Mrs.\ Sarah Clay (formerly Miss S.~Ellis)''.\\[4pt]
\multicolumn{3}{l}{\textit{The Diaries for 1826--1834 (not in Leybourn)}}\\
Mrs S.~Clay, of Moulton & at least nine questions in 1828, including Q1465 & Almost
  certainly Sarah Ellis, above.\\
Miss Jane Thoburn, of Blyth & Q1479 (1829) & First printed answer, ``aged 12 years''. A
  Mary Thoburn and a Miss Thoburn appear in the lists of the same years
  (\textit{The Gregory years}).\\
Miss L.~L. & Q1557 and at least seven others (1834) & Her solution to Q1557 is printed
  first. A footnote to Q1566 compares her with Agnesi and Mrs Somerville.\\[4pt]
\multicolumn{3}{l}{\textit{The Diaries for 1835--1840 (not in Leybourn)}}\\
Miss L.~L. & Q1599, 1601, 1606, 1607 (all 1837), probably 1604 and 1608 & Also
  answered two of the Diary's queries in 1837.\\
Mary Forster & Q1599 (1837), perhaps another with J.~Forster & \\
\bottomrule
\end{longtable}

\section*{What we know of them}

For almost all of these women the Diary is the only record. Leybourn gives a name and
sometimes a town, and nothing more. Three cases are worth setting out.

\subsection*{Miss Nancy Mason, of Clapham}

Her Clapham is the village under Ingleborough in the Yorkshire Dales, not Clapham in
London. Sylvia Harrop went looking for her in 2010.\footnote{S.~Harrop, ``The Enigma of
Miss Nancy Mason of Clapham, Eighteenth Century Mathematician Extraordinaire'',
\textit{North Craven Heritage Trust Journal} (2010). Online at
\url{http://settle.ddns.net/NCHTJ2010/2010/NancyMason.html}. Everything in this
subsection is from Harrop, except where noted.}

Nancy first appears in the Diary in 1792, aged about nineteen, and is there every year
until 1798. Then she disappears. Her question of 1795 asks for her own age on 1~January
of that year, and the answer is 21 years, 11 months and 7 days. That puts her birth
in late January 1773.\footnote{Our arithmetic, not Harrop's.} She was not only a
mathematician. She wrote verse charades and long answers to the Diary's ``queries'', on
whether a virtuous old maid or a young widow would make the more charming wife, and on
the origin of valentines. In 1793 she shared a prize with the Reverend Mr Ewbank, six
Diaries each. In 1796 Ewbank printed a poem to her, offering, half in jest, to ``tie the
knot'' if he visited. A modern reader, Mike Slater, judged her mathematics to be at
about today's A-level standard.

The parish registers of Clapham, Giggleswick and Bentham have no baptism or burial for
her, and the Quaker records have nothing either. Harrop concludes that she was not from
the district. Her guess is that Nancy was staying with relatives or friends of her
family, probably a well-off or clerical one. A Mason family lived at Newby, close by,
and its son James was almost exactly her age. A ``Mr James Mason of Clapham'' answered
one Diary question, in 1800. Harrop considers whether ``Nancy Mason'' was a pen-name and
rejects it: by the 1790s women had no need to hide.

Why did she stop? Harrop's likeliest answer is marriage, at about twenty-five. The Diary
itself had remarked in 1792 on three ``fair and ingenious sisters'' who stopped writing
one by one as each married. It may be the commonest ending in this whole dossier. We
simply cannot see it.

\subsection*{Mary Wright and Mary Nelson}

Albree and Brown print a suggestion that Mrs Mary Nelson of Newmarket, who set the
prize-money question in 1719, was Mary Wright under her married name.\footnote{J.~Albree and S.~H.~Brown, ``A valuable monument of mathematical genius'', \textit{Historia Mathematica} 36 (2009), 10--47, p.~17.}
The suggestion was Shelley Costa's, made tentatively in her doctoral dissertation (Cornell, 2000).
She has since found evidence that disproves it, and has told us so: they were
two different women.\footnote{S.~Costa, personal communication, October 2026.} We
repeat the point here only because the published version is wrong.

Mary and Anna Wright were sisters: in 1711 Tipper printed Mary's answer to the prize
question ``and likewise of the Enigma by her Sister Mrs.\ Anna''.\footnote{\textit{The
Ladies' Diary} for 1711, p.~31, reproduced in S.~A.~Costa, \textit{The Ladies' Diary:
Society, Gender and Mathematics in England, 1704--1754} (PhD dissertation, Cornell, 2000),
Figure~3.2.} Costa's dissertation goes further, tracing them and Barbara Sidway in
parish and county records. That part of her work is unpublished, and we do not draw on
it here.

\subsection*{Elizabeth Beighton}

She set and answered nothing that we can find, but for sixteen years the Diary was hers.
Even before Henry Beighton died, Charles Hutton wrote in 1775, the mathematical part was
composed by his friend Anthony Thacker, while ``the enigmatical parts were managed by his
amiable wife''. When Henry died in 1743, Elizabeth carried the Diary on as his widow, and
the Stationers' Company recognised her as its ``author'' throughout her widowhood, until
her death in 1759. She edited the enigmas herself. For the mathematics, and for the
astronomy an almanac needs, she relied on others: first Thacker, who died in 1744, then
Robert Heath, who
ran the mathematics from 1745 until the Company dismissed him in 1753, then Thomas
Simpson. The histories written afterwards, from Hutton's own onwards, gave the dates of
the men and left her out, and so made her part look marginal.\footnote{S.~A.~Costa,
\textit{The Ladies' Diary: Society, Gender and Mathematics in England, 1704--1754}, PhD
dissertation, Cornell University (2000), pp.~324--327, quoting C.~Hutton, \textit{The
Diarian Miscellany} (London, 1775), vol.~I, preface, p.~v, and H.~Ellis (ed.),
\textit{Letters of Eminent Literary Men} (London, 1843), p.~304. Thacker's death in 1744 is from R.~and P.~Wallis (eds), \textit{Index of British
Mathematicians, Part III, 1701--1800} (Newcastle upon Tyne, 1993), p.~135, cited by Costa
(p.~327, n.~25). The points on the
enigmas and on her reliance on Thacker and Heath are confirmed by Costa, personal
communication, October 2026. J.~Albree and S.~H.~Brown, ``A valuable monument of mathematical genius'', \textit{Historia Mathematica} 36 (2009), Table~1, list her as an editor for 1744 only.}
So a woman ran the Diary through the very years when, as the table above shows, women had
stopped setting its questions.

\subsection*{The rest}

For the others we found nothing outside the Diary. ProofWiki says of Barbara Sidway
only that ``little is known'' about her. The published studies (Perl, Costa's \textit{Osiris} article, Albree and
Brown, Swetz) treat the women as a group. The one sustained attempt to trace them in
the records is Costa's dissertation, mentioned above, which is unpublished.
The towns the Diary gives are the place to begin, as Harrop began with Clapham:

\begin{center}
\begin{tabular}{l l}
\toprule
Name & Town as printed\\
\midrule
Mrs Mary Nelson & Newmarket\\
Xantippe & Helmsley\\
Miss Betty Claxton & Benwell\\
Miss Maria Middleton & Eden, near Durham\\
Miss M.~Castieau & Shropshire\\
Miss Diana Browne & Honiton (1780)\\
Mrs Macdonald & Spindlestone, near Bamburgh\\
Miss Mary Weston & Gainsborough\\
Miss Susan Jackson & Mile End\\
Mrs Elizabeth Bryhon & Helmsley\\
\bottomrule
\end{tabular}
\end{center}

Castieau and Bryhon are rare surnames. The parish registers for those places may well
find them.

\section*{Pen-names, and the case of Ann Nicholls}

``Miss Ann Nicholls'' (also printed Nichols and Anna Nicholls) set ten questions between
1760 and 1768: 462, 475, 487, 516, 527, 543, 551, 556, 567 and 586. Most are astronomy:
sunrise at two towns, a lunar eclipse, the Pleiades and Aldebaran in one vertical
circle. Leybourn's index lists her as ``Nichols, Ann. (Mr.\ W.~Wales)''. William Wales
(c.\ 1734--1798) was the astronomer on Cook's second voyage. He later became master of
the Royal Mathematical School at Christ's Hospital. The index gives him several other
pseudonyms: Nauticus, Geometricus, Gemini and G.~Cetii.

Leybourn knew many of these contributors, and his identifications elsewhere check out.
He names Thomas Simpson as Hurlothrumbo, Kubernetes and Patrick O'Cavanah, and Landen as
Peter Puzzlem. So his word should be taken. The ten ``Ann Nicholls'' questions are a
man's, written under a woman's name.

It cuts both ways. The female pen-names of the early volumes, Xantippe, Adrastea,
Silvia, Hesychia and Anna Philomathes, may also hide men. Leybourn identifies none of
them, and the early Diary had a real community of women contributors to draw on, so
most are probably genuine. But the count of women setters is best read as an upper
bound for pen-names and a lower bound for women overall. Some women surely signed with
initials or men's names.

One pen-name is missing from Leybourn altogether. Albree and Brown note that Beighton
answered a ``Blowzabella'' as a woman in 1726, and that she won a prize for one of her
mathematical answers in 1740.\footnote{J.~Albree and S.~H.~Brown, ``A valuable monument of mathematical genius'', \textit{Historia Mathematica} 36 (2009), 10--47, p.~17, citing the Diary for 1726,
p.~20, and 1740, p.~40.} She is not in Leybourn's index, which supports the second
caution above: the index undercounts the women solvers.

\section*{The pattern}

\begin{itemize}
\item \textbf{1709--1724.} Women both set and answer. Barbara Sidway, Mary Nelson and the
two Wrights appear repeatedly. Several questions are about the women themselves: their
ages, their fortunes, their suitors.
\item \textbf{1725--1748.} No question in Leybourn's first volume is set by a woman after
Adrastea's prize question of 1724. The mathematics grows harder: fluxions, curves,
mechanics. The named contributors are now mostly schoolmasters, surveyors and excisemen.
\item \textbf{1748--1776.} A short revival as the second volume opens (Amanda, Miss
Munlove, Rosamond), then Miss Atkinson, Miss T.~S---e and Miss S.~T. Their questions are
mostly arithmetic in verse.
\item \textbf{1777--1802.} Miss Nancy Mason and Miss Sarah Cowen set questions. Women
solvers return in the 1790s: Miss Betty Claxton, Miss Maria Middleton, Miss Castieau.
\item \textbf{1802--1816.} One setter, Mrs Macdonald. But Miss Mary Weston answers eight
questions, more than any woman since 1720.
\item \textbf{After Leybourn, 1817--1840.} Albree and Brown know of only three women who
set mathematical questions between 1817 and 1836, and none after. In the last issue,
for 1840, no woman appears in the mathematics at all. Two newly found volumes add names. The issues for 1816--1825 have Mary Knight and
Susan Jackson setting questions, and five or six women answering, among them Sarah Ellis,
who went on answering after her marriage, as Mrs Sarah Clay.\footnote{\textit{The
Ladies' Diary} for 1816--1825, digitised by Google Books.} The issues for 1826--1834 add
Miss Jane Thoburn of Blyth, who in 1829, ``aged 12 years'', has the first printed answer
to a question, and show Mrs Clay still answering in 1828. By 1834 Miss L.~L. is answering
in quantity. Her solution to Question 1557 is printed first, and against her name
in the list of solvers of Question 1566 the editor put a footnote: ``We admire
the talents of our fair correspondent. If she will only make up her mind to carry her
own solutions through, instead of referring to Simpson, Hymer, \&c.\ she may soon vie
with Donna Agnesi, and our own Mrs.\ Somerville.''\footnote{\textit{The Ladies' Diary} for
1825--1834, digitised by Google Books; the note is in the issue for 1834.} The bound volume of the last six
issues at Harvard adds another: a Miss L.~L., who set questions in 1834 and 1836 and was
among the most regular solvers of 1837. She is probably one of Albree and Brown's
three, though they mention her only once, as one of seven solvers credited for a single
answer. Mary Forster answered in 1837. In the enigmas, rebuses and verse of the same
issues, women are everywhere.\footnote{\textit{The Ladies' Diary} for 1835--1840,
Harvard College Library copy, digitised by HathiTrust,
\url{https://hdl.handle.net/2027/hvd.hxvmnt}; Albree and Brown, p.~28, n.~27.}
\end{itemize}

\section*{Further reading}

J.~Albree and S.~H.~Brown, ``A valuable monument of mathematical genius'': \textit{The
Ladies' Diary} (1704--1840), \textit{Historia Mathematica} 36 (2009).

T.~Perl, The \textit{Ladies' Diary} or Woman's Almanack, 1704--1841, \textit{Historia
Mathematica} 6 (1979).

S.~Costa, The \textit{Ladies' Diary}: Gender, Mathematics, and Civil Society in
Early-Eighteenth-Century England, \textit{Osiris} 17 (2002).

S.~A.~Costa, \textit{The Ladies' Diary: Society, Gender and Mathematics in England,
1704--1754}, PhD dissertation, Cornell University (2000).

F.~J.~Swetz, \textit{The Impact and Legacy of The Ladies' Diary (1704--1840): A Woman's
Declaration} (MAA Press, 2021), especially the chapter ``Women and the Diary''.

S.~Harrop, ``The Enigma of Miss Nancy Mason of Clapham'', \textit{North Craven Heritage
Trust Journal} (2010).

\end{document}
